\section{The \emph{Celestial Hierarchy} of Pseudo-Dionysius}
\label{sec:celestial-hierarchy}

No single treatise outside Sacred Scripture shaped Catholic angelology as deeply as \emph{De caelesti hierarchia} (\emph{The Celestial Hierarchy}). Presenting himself as Dionysius the Areopagite, the Athenian judge converted by St.~Paul's preaching on the Areopagus (Acts~17:34), and addressing his companion ``Fellow Presbyter Timothy,'' the author composed his four treatises and ten letters around the close of the fifth or the opening of the sixth century in a Syrian Christian milieu steeped in the liturgy of the East and in conversation with the Neoplatonism of Proclus. Both facts stand together: modern historical criticism places the corpus around A.D.~500, while the Church's reception from the sixth century onward prized the work for the theological depth of its doctrine. St.~Gregory the Great cited Dionysius by name in Rome as ``an ancient and venerable Father'' (\emph{Homily~34 on the Gospels}~\S 12); St.~Maximus the Confessor commented upon his corpus and St.~John of Damascus adopted his nine choirs verbatim (\emph{Exposition of the Orthodox Faith} II.3); John Scotus Eriugena translated him into Latin in the ninth century, followed by Hugh of St.~Victor, St.~Albert the Great, and St.~Thomas Aquinas, who wove \emph{The Celestial Hierarchy} into the very framework of the \emph{Summa theologiae}.

One terminological note prevents a recurring confusion in English translations of the middle triad. In chapter~8, the Greek text names the three middle choirs \gk{Kyriot\={e}tes}, \gk{Dynameis}, and \gk{Exousiai}, rendered in the standard Latin tradition as \latin{Dominationes}, \latin{Virtutes}, and \latin{Potestates}---traditional English \textbf{Dominations} (or Dominions), \textbf{Virtues}, and \textbf{Powers}. John Parker's 1899 English translation, quoted throughout the chapter exposition below, renders those same three Greek nouns as \textbf{Lordships}, \textbf{Powers}, and \textbf{Authorities}. Thus Parker's ``Powers'' always corresponds to Latin \latin{Virtutes} (Virtues), while Parker's ``Authorities'' corresponds to Latin \latin{Potestates} (Powers).

\subsection{I: Divine Illumination and Sacred Veils (PG~3, cols.~119--124)}
\label{sec:ch-i}
\emph{To my Fellow Presbyter Timothy. Dionysius the Presbyter. That every divine illumination, whilst going forth lovingly to the objects of its forethought under various forms, remains simplex. Nor is this all. It also unifies the things illuminated.}

Opening with James~1:17 and Romans~11:36, Dionysius traces the double rhythm of divine light: proceeding outward from the Father (\gk{proodos}) under manifold forms, it yet stays undivided in itself and converts creatures upward (\gk{epistroph\={e}}) into unifying simplicity (\S 1). Invoking Jesus as the ``Paternal Light'' (\S 2), he explains why God clothes immaterial mysteries in sensible images and sacramental liturgy (\S 3). Because human intellect cannot leap straight to pure contemplation, divine providence wraps its ray in ``variegated sacred veils'' (\gk{poikilois hierois parapetasmasin})---scriptural figures, fragrances, lights, clerical orders, and the Eucharist---lifting the mind from visible beauties to the invisible splendors of heaven:
\begin{quotation}
For it is not possible that the supremely Divine Ray should otherwise illuminate us, except so far as it is enveloped, for the purpose of instruction, in variegated sacred veils, and arranged naturally and appropriately, for such as we are, by paternal forethought. (\CH{1.2})
\end{quotation}
St.~Thomas Aquinas adopts this chapter directly at the opening of the \emph{Summa} (\ST{I, q.~1, a.~9}) to justify Scripture's use of corporeal metaphors, and returns to it when grounding angelic hierarchy and illumination in the divine light (\ST{I, q.~108, a.~1}; \ST{I, q.~111, a.~1}).

\subsection{II: Dissimilar Symbols and Apophatic Theology (PG~3, cols.~136--146)}
\label{sec:ch-ii}
\emph{That Divine and Heavenly things are appropriately revealed, even through dissimilar symbols.}

Warning against crude literalism---angels are not really many-footed beasts, fiery wheels, or oxen (\S 1)---Dionysius defends Scripture's striking use of humble and incongruous imagery (\S\S 2--3). Because God surpasses every created essence, negations (\gk{apophaseis}) speak more truly of Him than affirmations (\gk{kataphaseis}). Luminous golden figures tempt earthly minds to mistake the image for the reality, whereas base images---lions, bears, even a worm (Ps~21:7 [22:6], \S 5)---jolt the soul upward and protect sacred mysteries from irreverence, while passions like anger or desire signify steadfast courage and inflexible divine love in pure spirits (\S 4):
\begin{quotation}
If, then, the negations respecting things Divine are true, but the affirmations are inharmonious, the revelation as regards things invisible, through dissimilar representations, is more appropriate to the hiddenness of things unutterable. (\CH{2.3})
\end{quotation}
Aquinas reproduces all three Dionysian reasons in \ST{I, q.~1, a.~9, ad~3}: dissimilar corporeal figures preserve minds from error, suit our apophatic knowledge of God in this life, and veil holy mysteries from the unworthy.

\subsection{III: The Definition and Purpose of Hierarchy (PG~3, cols.~164--168)}
\label{sec:ch-iii}
\emph{What is Hierarchy? and what the use of Hierarchy?}

Here Dionysius gives the classical definition of hierarchy (\gk{hierarchia}) that governs both heaven and the Church (\S 1):
\begin{quotation}
Hierarchy is, in my judgment, a sacred order and science and operation, assimilated, as far as attainable, to the likeness of God, and conducted to the illuminations granted to it from God, according to capacity, with a view to the Divine imitation. (\CH{3.1})
\end{quotation}
Its end (\S 2) is assimilation and union with God, making each member a spotless mirror and a ``fellow-worker with God'' (1~Cor~3:9), receiving divine radiance and passing it ungrudgingly downward through the three hierarchical acts of purification, illumination, and perfection (\gk{katharsis, ph\={o}tismos, telei\={o}sis}): ``For it is an Hierarchical regulation that some are purified and that others purify; that some are enlightened and others enlighten; that some are perfected and others perfect'' (\CH{3.2--3}). Aquinas cites this definition and its triadic activity across \ST{I, q.~108, aa.~1, 2, 4} and \ST{I, q.~106, aa.~1, 3} (as well as \emph{De veritate} q.~9, aa.~2--3), clarifying that among sinless angels ``purification'' removes simple lack of knowledge rather than moral stain.

\subsection{IV: The Appellation ``Angels'' and Mediated Revelation (PG~3, cols.~177--184)}
\label{sec:ch-iv}
\emph{What is meant by the appellation ``Angels?''}

Tracing the ladder of participation from inanimate beings (sharing Being), through living creatures (sharing Life), to rational intellects (sharing Wisdom, \S 1), Dionysius explains that the heavenly minds are pre-eminently named ``Angels'' (messengers) because divine illumination reaches them at first hand and passes through them to us, as when the Mosaic Law was delivered through angels (\S 2). Because no one has seen the hidden essence of God, Old Testament theophanies were mediated through celestial powers according to the law that ``through the first, the second are brought to the Divine Being'' (\S 3). In \S 4, the argument culminates in Christ: angels were first initiated into the Incarnation (Gabriel to Zachary and Mary, the angel to Joseph, the host to the shepherds), and the incarnate Lord Himself respected the hierarchical order He established:
\begin{quotation}
for I perceive that even Jesus Himself, the superessential Cause of the super-heavenly Beings, when He had come to our condition, without change, did not overstep the good order which becomes mankind, which Himself arranged and took, but readily subjected Himself to the dispositions of the Father and God, through Angels... or that even Jesus Himself, when He had come to manifest the good work of our beneficent salvation, was called Angel of Great Counsel. (\CH{4.4})
\end{quotation}
Aquinas relies on Chapter~IV to establish hierarchical illumination (\ST{I, q.~106, a.~3}), angelic enlightenment of men (\ST{I, q.~111, a.~1}), angelic knowledge of grace (\ST{I, q.~57, aa.~1, 5}), personal guardianship (\ST{I, q.~113, a.~2}), and the fittingness of Gabriel's Annunciation to Mary (\ST{III, q.~30, a.~2}).

\subsection{V: Why All Heavenly Beings Are Called Angels in Common (PG~3, col.~196)}
\label{sec:ch-v}
\emph{For what reason all the Heavenly Beings are called, in common, Angels.}

In a brief, incisive chapter, Dionysius solves a lexical puzzle: why does Scripture call every heavenly spirit an ``angel'' while assigning the name ``Angels'' specifically to the lowest of the nine orders? He formulates the principle of asymmetrical containment:
\begin{quotation}
Now, we affirm that throughout every sacred ordinance the superior ranks possess the illuminations and powers of their subordinates, but the lowest have not the same powers as those who are above them. (\CH{5})
\end{quotation}
Higher ranks like Seraphim or Thrones possess the messenger-power in abundance and may rightly be called angels, whereas the lowest rank lacks the supreme endowments of Principalities or Seraphim and never bears their names. Aquinas adopts this principle verbatim to explain the naming of all nine choirs (\ST{I, q.~108, a.~5}) and the assignment of individual guardianship to the lowest choir (\ST{I, q.~113, a.~3}).

\subsection{VI: The Three Threefold Orders of Heaven (PG~3, cols.~200--201)}
\label{sec:ch-vi}
\emph{Which is the first Order of the Heavenly Beings? which the middle? and which the last?}

Opening with sober apophatic humility, Dionysius declares that only God and the angels themselves know fully how many heavenly ranks exist and how they are ordered; human knowledge reaches only what Scripture reveals:
\begin{quotation}
How many, and of what sort, are the Orders of the supercelestial Beings, and how the Hierarchies are classified amongst themselves, I affirm, the deifying Author of their consecration alone distinctly knows... The Word of God has designated the whole Heavenly Beings as nine, by appellations, which shew their functions. These our Divine Initiator divides into three threefold Orders. (\CH{6})
\end{quotation}
Citing his teacher Hierotheus, he assigns three choirs to each hierarchy: the first (immediately around God) holds Thrones, Cherubim, and Seraphim; the middle holds Authorities, Lordships, and Powers; the third holds Angels, Archangels, and Principalities. Whereas Chapter~VI lists each triad's members in rhetorical sequence, Chapters~VII through~IX (confirmed explicitly at \CH{11.1}) fix the exact descending order of all nine ranks. Aquinas takes this triadic architecture directly into \ST{I, q.~108, aa.~1--3}, and cites Chapter~VI also on angelic self-knowledge (\ST{I, q.~56, a.~1}) and the fall (\ST{I, q.~63, a.~7}).

\subsection{VII: The First Hierarchy: Seraphim, Cherubim, and Thrones (PG~3, cols.~205--212)}
\label{sec:ch-vii}
\emph{Concerning the Seraphim and Cherubim and Thrones, and concerning their first Hierarchy.}

Expounding the supreme triad stationed immediately around God (\S 1), Dionysius interprets the Hebrew names: \textbf{Seraphim} means ``kindling or burning,'' signifying perpetual motion around divine things, purifying heat, and unquenchable ardor; \textbf{Cherubim} means ``a fulness of knowledge or stream of wisdom,'' signifying unimpeded contemplation and generous imparting of wisdom; \textbf{Thrones} signifies passionless exaltation above every earthly baseness and God-bearing receptivity (\gk{theophoria}). These supreme minds are purified not from moral stain, but by exaltation above all remissness, illuminated directly by immediate communion with God, and perfected without intermediary (\S 2). From Psalm~23:7--10 [24:7--10] (lower angels asking higher angels, ``Who is this King of Glory?'') and Isaias~63:1--2 (even the highest spirits questioning Jesus on His reddened garments), Dionysius demonstrates both inter-angelic instruction and Christ's direct illumination of the highest choir (\S 3):
\begin{quotation}
But to sum up, I may say this, not inappropriately, that the reception of the supremely Divine Science is, both purification, and enlightenment, and perfecting,---purifying, as it were, from ignorance, by the knowledge of the more perfect revelations imparted to it according to fitness, and enlightening by the self-same Divine knowledge... and perfecting further, by the self-same Light, through the abiding science of the mysteries made clearly manifest. (\CH{7.3})
\end{quotation}
Section~4 contemplates their perpetual choral dance (\gk{choreia}) around the tri-subsistent Monad, echoing the Trisagion of Isaias~6:3 and Ezechiel~3:12. Aquinas draws on Chapter~VII throughout the \emph{Summa}: on mysteries of grace (\ST{I, q.~57, a.~5}), the fall of the highest angels (\ST{I, q.~63, a.~7}), illumination and speech of lower angels to higher (\ST{I, q.~106, aa.~1--2, 4}; \ST{I, q.~107, a.~2}), and the names and ranks of the first hierarchy (\ST{I, q.~108, aa.~1, 4--6}; \ST{I, q.~112, a.~4}).

\subsection{VIII: The Middle Hierarchy: Dominations, Virtues, and Powers (PG~3, cols.~237--242)}
\label{sec:ch-viii}
\emph{Concerning Lordships and Powers and Authorities, and concerning their middle Hierarchy.}

Turning to the middle triad---Greek \gk{Kyriot\={e}tes, Dynameis, Exousiai}, Latin \latin{Dominationes, Virtutes, Potestates}---Dionysius expounds their cosmic governance (\S 1):
\begin{quotation}
I think, then, that the explanatory name of the Holy Lordships denotes a certain unslavish elevation, free from all grovelling subserviency, as becomes the free... and the appellation of the Holy Powers denotes a certain courageous and unflinching virility, for all those Godlike energies within them... and that the appellation of the Holy Authorities, of the same rank as the Divine Lordships and Powers, (denotes) the beautiful and unconfused good order, with regard to the Divine receptions, and the discipline of the supermundane and intellectual authority. (\CH{8.1})
\end{quotation}
Purified and enlightened at second hand through the first triad, the middle triad demonstrates the universal law of mediation (\S 2): ``For, this is divinely fixed absolutely by the Divine source of order that, through the first, the second partake of the supremely Divine illuminations.'' Dionysius corroborates inter-angelic command from Zacharias~2:3--4, Ezechiel~9--10 (the Cherub casting fire into the hands of the linen-vested angel), and Daniel~8:16 (a higher voice commanding Gabriel). Aquinas builds on Chapter~VIII to prove that higher angels enlighten lower (\ST{I, q.~106, a.~1}), to define the middle triad (\ST{I, q.~108, aa.~5--6}), and to determine that Dominations direct external missions without being sent to execute them (\ST{I, q.~112, a.~4}).

\subsection{IX: The Last Hierarchy: Principalities, Archangels, and Angels (PG~3, cols.~257--264)}
\label{sec:ch-ix}
\emph{Concerning the Principalities, Archangels, and Angels, and concerning their last Hierarchy.}

The third triad presides directly over human hierarchies and nations (\S\S 1--2). \textbf{Principalities} (\gk{Archai}) turn wholly toward the ``super-princely Prince'' and lead in princely order; \textbf{Archangels} stand in the middle, sharing upward with Principalities in marshalling unity and downward with Angels in announcing illumination to men; and \textbf{Angels} complete the celestial series, ``more properly named Angels by us than those of higher degree, because their Hierarchy is occupied with the more manifest, and is more particularly concerned with the things of the world'' (\CH{9.2}). Citing Daniel~10:13, 21 (Michael as ruler of Israel) and Deuteronomy~32:8 LXX, Dionysius faces a classic objection (\S\S 3--4): if holy angels guided every nation, why did only Israel worship the true God? He exonerates the national angels completely---Gentiles strayed through their own self-will, while righteous non-Israelites like Melchizedek, and even Pharaoh through Joseph and Nebuchadnezzar through Daniel, received true light---and bars every polytheistic reading of Israel's election:
\begin{quotation}
For there is one Prince and Providence over all. And never must we think that the Godhead is leader of Jews by lot, and that Angels, independently, or as of equal rank, or in opposition, or that certain other gods, preside over the other nations... demonstrating distinctly that there is one Providence of the whole, superessentially established above all the powers, unseen and seen, and that all the Angels who preside over each nation, elevate, as far as possible, those who follow them with a willing mind, to It as their proper Head. (\CH{9.4})
\end{quotation}
Aquinas draws upon Chapter~IX for the third hierarchy (\ST{I, q.~108, aa.~5--6}) and for the assignment of individual guardianship to the lowest order while Principalities or Archangels watch over kingdoms (\ST{I, q.~113, a.~3}).

\subsection{X: Summary of the Angelic Discipline (PG~3, cols.~272--274)}
\label{sec:ch-x}
\emph{A Repetition and Summary of the Angelic discipline.}

Summarizing the cascade from God through the three triads to our human hierarchy (\S 1), Dionysius calls every angel an ``interpreter'' (\gk{herm\={e}neus}) of those above it (\S 2) and reveals that the triadic law repeats fractally within every single choir---proving from Isaias~6:3 (``the most Divine Seraphim cry one to another'') that higher Seraphim impart divine knowledge to lower Seraphim---and even inside every individual created intellect (\S 3):
\begin{quotation}
I might add this not inappropriately, that each heavenly and human mind has within itself its own special first, and middle, and last ranks, and powers, manifested severally in due degree... For there is nothing that is self-perfect, or absolutely without need of perfecting, except the really Self-perfect and preeminently Perfect. (\CH{10.3})
\end{quotation}
Aquinas turns Chapter~X into a striking metaphysical authority: because Aristotle shows there is no ``first and second'' within one species, the presence of first, middle, and last ranks even inside a single choir supports the Thomist thesis that every angel is its own distinct species (\ST{I, q.~50, a.~4}); he also invokes it on hierarchical harmony and guardianship (\ST{I, q.~108, a.~3}; \ST{I, q.~113, aa.~2--3}).

\subsection{XI: Why All Heavenly Beings Are Called Heavenly Powers (PG~3, cols.~284--285)}
\label{sec:ch-xi}
\emph{For what reason all the Heavenly Beings, in common, are called Heavenly Powers.}

Raising a sharp objection against his own rule in Chapter~V, Dionysius asks why Scripture calls all celestial minds ``Heavenly Powers'' (\gk{ouraniai dynameis}) when Angels, Archangels, Principalities, and Authorities stand \emph{below} the choir of Powers (\gk{Dynameis} / Virtues)---incidentally locking in his descending rank order (\S 1). He answers from the threefold metaphysical structure of every created spirit (\S 2):
\begin{quotation}
But, inasmuch as all the Divine Minds, by the supermundane description given of them, are distributed into three,---into essence, and power, and energy,---when we speak of them all, or some of them, indiscriminately, as Heavenly Beings or Heavenly Powers, we must consider that we manifest those about whom we speak in a general way, from their essence or power severally. (\CH{11.2})
\end{quotation}
St.~Thomas Aquinas adopts this triad of essence, power, and operation (\gk{ousia, dynamis, energeia}) directly in the \emph{sed contra} of \ST{I, q.~54, a.~3} to establish the real distinction between an angel's essence and its operative power of intellect.

\subsection{XII: Why Human Hierarchs Are Called Angels (PG~3, cols.~292--293)}
\label{sec:ch-xii}
\emph{Why the Hierarchs amongst men are called Angels.}

Confronting Malachy~2:7, which calls the human priest ``Angel of the Sovereign Lord'' (\S 1), Dionysius complements Chapter~V with the rule of subordinate participation (\S 2):
\begin{quotation}
for, as the first possess abundantly the saintly characteristics of the inferior, so the last possess those of the superior, not indeed in the same degree, but subordinately. There is, then, as I think, nothing absurd, if the Word of God calls our Hierarch, Angel, since he participates, according to his own capacity, in the messenger characteristic of the Angels. (\CH{12.2})
\end{quotation}
He confirms this analogical naming from Psalm~81:6 [82:6], where Scripture even calls angels and holy men ``gods'' by deifying participation (\S 3). Aquinas cites Chapter~XII in \ST{I, q.~55, a.~3} (higher angels understand through more universal species, shared partially by lower ranks) and \ST{I, q.~106, a.~4}.

\subsection{XIII: Isaias's Purification by the Seraph (PG~3, cols.~300--308)}
\label{sec:ch-xiii}
\emph{For what reason the Prophet Isaiah is said to have been purified by the Seraphim.}

Here Dionysius meets the stiffest scriptural challenge to strict hierarchical mediation: in Isaias~6:6--7, one of the \emph{Seraphim}---the highest choir of all---flies down with a live coal to cleanse the prophet's lips (\S 1). After noting a simple first answer---that an inferior angel was called ``seraph'' homonymously because his mission worked by fire (\S 2)---Dionysius develops a deeper explanation through the analogies of sunlight traversing transparent media and fire warming water through a vessel (\S 3). Because every heavenly illumination flows from God as Cause through the supreme triad downward, the lower angel who actually ministered to Isaias ascribed his fiery purification upward (\S 4):
\begin{quotation}
just as our Hierarch, when purifying or enlightening through his Leitourgoi or Priests, is said himself to purify and enlighten, since the Orders consecrated through him attribute to him their own proper sacred operations; so also the Angel who effected the purification of the Theologian attributes his own purifying science and power to God, indeed, as Cause, but to the Seraphim as first-operating Hierarch. (\CH{13.4})
\end{quotation}
Dionysius closes with quiet modesty, inviting Timothy to embrace whichever account seems better. St.~Thomas Aquinas examines this exact question in \ST{I, q.~112, a.~2} beside St.~Gregory the Great (\emph{Homily~34}~\S\S 12--13) and sides firmly with Dionysius: because the order of grace has no higher order for whose sake it would be bypassed, the highest angels are never sent to external ministry, and the ministering angel in Isaias~6 bore the name of Seraph by virtue of the purifying mission received from above. He also invokes Chapter~XIII alongside Chapter~IV when grounding individual guardianship in hierarchical mediation (\ST{I, q.~113, a.~2}).

\subsection{XIV: The Traditional Number of the Angels (PG~3, cols.~321--322)}
\label{sec:ch-xiv}
\emph{What the traditional number of the Angels signifies.}

Expounding Daniel~7:10 (``thousand thousands, and myriad myriads''), Dionysius shows that Scripture heaps up the highest limits of human numeration to teach that angelic multitude exceeds every earthly reckoning, while remaining finite and definitely known to the angels' own intellects and to the omniscient Creator, the ``encompassing Term'' of all creatures:
\begin{quotation}
For many are the blessed hosts of the supermundane minds, surpassing the weak and contracted measurement of our material number, and being definitely known by their own supermundane and heavenly intelligence and science alone, which is given to them in profusion by the supremely Divine and Omniscient Framer of Wisdom. (\CH{14})
\end{quotation}
Aquinas cites Chapter~XIV directly in \ST{I, q.~50, a.~3} to reject Aristotle's restriction of separate substances to the number of celestial spheres and to affirm that the angels surpass all material multitude, returning to it again in \ST{I, q.~112, a.~4}.

\subsection{XV: Symbolic Figures of the Angelic Powers (PG~3, cols.~325--340)}
\label{sec:ch-xv}
\emph{What are the morphic likenesses of the Angelic Powers? what the fiery? what the anthromorphic [sic]? what are the eyes? what the nostrils? what the ears? what the mouths? what the touch? what the eyelids? what the eyebrows? what the prime? what the teeth? what the shoulders? what the elbows and the hands? what the heart? what the breasts? what the back? what the feet? what the wings? what the nakedness? what the robe? what the shining raiment? what the sacerdotal? what the girdles? what the rods? what the spears? what the battle-axes? what the measuring lines? what the winds? what the clouds? what the brass? what the electron? what the choirs? what the clapping of hands? what the colours of different stones? what the appearance of the lion? what the appearance of the ox? what the appearance of the eagle? what the horses? what the varieties of coloured horses? what the rivers? what the chariots? what the wheels? what the so-called joy of the Angels?}

The closing chapter forms a complete symbolic lexicon of biblical angelology (\S\S 1--9). After noting that the same orders are shown both ruling and being ruled relative to lower and higher ranks (\S 1), Dionysius singles out \textbf{fire} as the highest scriptural image of angelic energy---ubiquitous, unmingled, luminous, elevating, and undiminished when shared (\S 2): ``I think, then, the similitude of fire denotes the likeness of the Heavenly Minds to God in the highest degree.'' Section~3 decodes the human form and every bodily feature: uprightness signifies intellect, eyes transparent receptivity to light, teeth the dividing of unified wisdom for inferiors, wings elevating lightness above earthly weight, and unshod feet unfettered simplicity. Subsequent sections interpret glowing robes and sacerdotal vestments (\S 4), rods, spears, battle-axes, and measuring lines (\S 5), swift winds and light-bearing clouds (\S 6), electron, brass, multicolored stones, and the indomitable lion (\S 7), and the strong ox, soaring eagle, and docile colored horses (\S 8). Section~9 expounds fiery rivers, chariots, the winged wheels named in Hebrew ``Gel, Gel'' (\gk{galgal})---``This shews, according to the Hebrew tongue, revolutions and revelations''---and the angels' joy over repentant sinners (Luke~15:10) as godlike, ungrudging gladness, ending in apophatic reverence:
\begin{quotation}
Other things, however, as being parallel to the things said, we have omitted, out of regard to the symmetry of the discourse; and the hiddenness, beyond our capacity, we have honoured by silence. (\CH{15.9})
\end{quotation}
Aquinas quotes Chapter~XV on the downward communication of unified light (\ST{I, q.~106, aa.~1, 4}), angelic speech (\ST{I, q.~107, a.~5}), and the figurative meaning of bodily limbs ascribed to angels in Scripture (\ST{I, q.~51, a.~3, ad~4}).

\subsection{The Dionysian and Gregorian Orderings of the Nine Choirs}
\label{subsec:dionysius-gregory-orderings}

Although Pseudo-Dionysius and St.~Gregory the Great agree that Sacred Scripture reveals nine angelic orders, they do not arrange them identically---indeed, Gregory himself prints two distinct sequences in his \emph{Homily~34 on the Gospels} (\S\S 7, 10) and his earlier \emph{Moralia in Job} (XXXII.23.48). The table below sets all three verified sequences side by side from highest to lowest.

{\small
\begin{longtable}{>{\raggedright\arraybackslash}p{0.07\linewidth}>{\raggedright\arraybackslash}p{0.28\linewidth}>{\raggedright\arraybackslash}p{0.28\linewidth}>{\raggedright\arraybackslash}p{0.31\linewidth}}
\toprule
\textbf{Rank} & \textbf{Pseudo-Dionysius} (\CH{7--9, 11.1}) & \textbf{St.~Gregory I} (\emph{Hom.~34.7, 10}) & \textbf{St.~Gregory I} (\emph{Moralia} XXXII.23.48) \\
\midrule
\endfirsthead
\toprule
\textbf{Rank} & \textbf{Pseudo-Dionysius} (\CH{7--9, 11.1}) & \textbf{St.~Gregory I} (\emph{Hom.~34.7, 10}) & \textbf{St.~Gregory I} (\emph{Moralia} XXXII.23.48) \\
\midrule
\endhead
\bottomrule
\endfoot
1 & Seraphim (\gk{Seraphim}) & Seraphim & Seraphim \\
2 & Cherubim (\gk{Cheroubim}) & Cherubim & Cherubim \\
3 & Thrones (\gk{Thronoi}) & Thrones & Powers (\latin{Potestates}) \\
4 & Dominations (\gk{Kyriot\={e}tes}) & Dominations (\latin{Dominationes}) & Principalities (\latin{Principatus}) \\
5 & Virtues (\gk{Dynameis}) & Principalities (\latin{Principatus}) & Virtues (\latin{Virtutes}) \\
6 & Powers (\gk{Exousiai}) & Powers (\latin{Potestates}) & Dominations (\latin{Dominationes}) \\
7 & Principalities (\gk{Archai}) & Virtues (\latin{Virtutes}) & Thrones (\latin{Throni}) \\
8 & Archangels (\gk{Archangeloi}) & Archangels (\latin{Archangeli}) & Archangels (\latin{Archangeli}) \\
9 & Angels (\gk{Angeloi}) & Angels (\latin{Angeli}) & Angels (\latin{Angeli}) \\
\end{longtable}
}

In \ST{I, q.~108, a.~6} (echoing his Pauline commentaries on Ephesians~1:21 and Colossians~1:16), St.~Thomas Aquinas compares Dionysius with Gregory's mature ordering in \emph{Homily~34} and shows why both are thoroughly reasonable without forcing them into artificial sameness. Both Doctors agree on seven of the nine choirs, differing only over \textbf{Virtues} and \textbf{Principalities}: Dionysius places Virtues between Dominations and Powers, and Principalities between Powers and Archangels, matching St.~Paul's ascending sequence in Ephesians~1:21 (``Above all principality and power and virtue and dominion''); Gregory places Principalities between Dominations and Powers, and Virtues between Powers and Archangels, matching Colossians~1:16 descending (``thrones, or dominations, or principalities, or powers'').

Aquinas shows that each Doctor follows a distinct, coherent principle of order. \textbf{Dionysius} derives the triads from the internal structure of contemplation and universal action: the first hierarchy considers the ultimate End in God Himself (Thrones receiving Him, Cherubim knowing His secrets, Seraphim united to Him in fire); the second governs universal disposition (Dominations appointing what is to be done, Virtues imparting strength for execution, Powers ordering how subordinates carry it out); the third executes particular ministries (Principalities leading nations and kingdoms, Archangels announcing supreme mysteries, Angels announcing particular messages to individual men). \textbf{Gregory}, on the other hand, arranges the orders beneath Dominations according to the hierarchy of subjects governed, following St.~Augustine's rule that lower bodies are ruled by spiritual creatures and evil spirits by good spirits: first come Principalities presiding over good spirits, then Powers coercing evil spirits, then Virtues exercising miracle-working sway over corporeal nature, and finally Archangels and Angels announcing great and lesser messages to men. By laying both rationales side by side, Aquinas honors Gregory's pastoral ordering while adopting Dionysius's triadic architecture as the systematic spine of the \emph{Summa}'s treatise on the angels.
